\documentclass[12pt]{article}

% ==== PAQUETES ====
\usepackage[utf8]{inputenc}      % Codificación
\usepackage[T1]{fontenc}         % Acentos y símbolos en PDF
\usepackage[spanish]{babel}      % Idioma
\usepackage{amsmath, amssymb}    % Símbolos matemáticos
\usepackage{geometry}            % Márgenes
\usepackage{lmodern}             % Fuente mejorada
\usepackage{enumitem}            % Mejoras en listas
\usepackage{hyperref}            % Enlaces
\geometry{margin=2.5cm}

% ==== DOCUMENTO ====
\begin{document}

\section*{Grados de libertad del test \textit{F} de White}

El test de White para heterocedasticidad se basa en una regresión auxiliar de los residuos al cuadrado $\hat{u}_i^2$ sobre todas las variables explicativas del modelo original, sus cuadrados y sus productos cruzados. Esta regresión auxiliar tiene como propósito detectar cualquier relación no lineal entre los regresores y la varianza de los errores, sin imponer una forma funcional específica.

\[
\hat{u}_i^2 = \delta_0 + \delta_1 x_{1i} + \dots + \delta_k x_{ki} + \delta_{k+1} x_{1i}^2 + \dots + \delta_{2k} x_{ki}^2 + \text{(productos cruzados)} + v_i
\]

El estadístico \textit{F} se construye para probar la hipótesis nula:

\[
H_0: \delta_1 = \delta_2 = \dots = \delta_q = 0 \quad \text{(homocedasticidad)}
\]

El estadístico \textit{F} tiene la forma:

\[
F = \frac{(R_{aux}^2 / q)}{(1 - R_{aux}^2) / (n - q - 1)} \quad \sim \quad F_{q,\;n - q - 1}
\]

donde:

\begin{itemize}
  \item $n$ es el número de observaciones,
  \item $q$ es el número de regresores en la regresión auxiliar (excluyendo el intercepto),
  \item $R_{aux}^2$ es el coeficiente de determinación de la regresión auxiliar.
\end{itemize}

El número total de variables en la regresión auxiliar (es decir, los grados de libertad del numerador del estadístico \textit{F}) es:

\[
q = k + \frac{k(k+1)}{2}
\]

donde $k$ es el número de regresores en el modelo original (sin incluir la constante). Esta fórmula incluye:

\begin{itemize}
  \item $k$ regresores originales,
  \item $k$ términos cuadráticos,
  \item $\frac{k(k-1)}{2}$ productos cruzados entre regresores distintos.

  Para determinar cuántos productos cruzados distintos se incluyen en la regresión auxiliar, consideramos todas las combinaciones posibles de pares de regresores distintos del modelo original. Esto se calcula usando combinatoria:

\[
\text{Número de productos cruzados distintos} = \binom{k}{2} = \frac{k(k - 1)}{2}
\]

Esta expresión representa el número de formas de seleccionar dos regresores diferentes de un total de $k$, sin importar el orden (es decir, $x_1x_2$ y $x_2x_1$ se consideran el mismo término).

\end{itemize}



\end{document}
